% !TeX root = ../main.tex

\section{Experimental Setup}
\label{sec:experimental-setup}

The experiments separate three questions: whether the descriptors rank
poisoned images above clean images, whether a threshold calibrated on earlier
tasks remains useful on a later task, and whether historical inversion
samples are sufficient for detection without labelled poisoning examples.
Feature-removal experiments address which parts of the supervised
representation contribute to these decisions. Paired training controls
measure the forgetting induced by the attack independently of the detector.

The completed main study uses CIFAR-100 with five continual-learning
implementations and three seeds. A smaller extension retains the full
miniImageNet and tinyImageNet datasets but restricts the learner to EWC.
The latter experiments were still running when this chapter was prepared;
their configuration is recorded separately below, without treating planned
runs as completed observations.

\subsection{Datasets and Task Sequences}
\label{subsec:setup-datasets}

CIFAR-100 contains 100 classes of $32\times32$ colour images, with 500
training and 100 test images per class~\cite{krizhevsky2009learning}.
The main study uses ten tasks of ten classes. The class order is fixed:
task $t$ contains global class identifiers $10t,\ldots,10t+9$. Seeds change
model training and attack generation, not the class order or the underlying
image collection.

The additional benchmarks follow the task counts used in the BrainWash
study~\cite{abbasi2024brainwash}: ten tasks for miniImageNet and twenty for
tinyImageNet. Both have ten classes per task. Table~\ref{tab:setup-datasets}
summarises the data used by the registered protocols. Images retain their
native benchmark resolution and are represented as floating-point RGB
arrays in $[0,1]$. Poisoned images are clipped back to this interval.

\begin{table}[tbp]
\centering
\small
\begin{tabular}{lrrrr}
\toprule
Dataset & Classes & Resolution & Tasks & Train/test per class\\
\midrule
CIFAR-100 & 100 & $32\times32$ & 10 & 500 / 100\\
miniImageNet & 100 & $84\times84$ & 10 & 500 / 100\\
tinyImageNet & 200 & $64\times64$ & 20 & 500 / 50\\
\bottomrule
\end{tabular}
\caption{Full benchmark splits. For tinyImageNet, ``test'' refers to the
labelled validation images; the unlabelled test split is not used.}
\label{tab:setup-datasets}
\end{table}

The prepared miniImageNet arrays preserve the existing training/test split
provided with the experimental assets. The tinyImageNet conversion keeps
all official training images and uses the annotated validation images for
evaluation of the continual learner. Its global class identifiers are
assigned by sorted class names. File hashes, class counts, image shapes
and the conversion record are checked before plan registration. No
class-balanced subset or reduced-image training run is substituted for the
full benchmark.

Two uses of the word ``test'' must be distinguished. The benchmark test
split measures the continual learner's class accuracy. The detector test
split is held out from a task's \emph{training} originals, because the
detector is intended to inspect incoming training data. These are different
image pools serving different measurements.

\subsection{Task Roles and Data Separation}
\label{subsec:setup-task-roles}

The source is incoming Task 1. The model used to describe it has learned
clean Task 0, but not Task 1. Source clean/BrainWash pairs fit the scaler
and logistic regression. Clean validation and reserve images from Tasks 1
and 4 calibrate the thresholds. The Task 4 detector test split is a
development diagnostic. Final CIFAR-100 evaluation uses incoming Task 9,
whose backbone has learned clean Tasks 0--8.

Each stage has its own clean pre-update checkpoint and historical
inversions. A source attack is not propagated through the training history
to create the target checkpoint. Although source and target classes are
disjoint, the later clean backbone has legitimately learned the source
task. The training histories should therefore not be described as
independent populations.

Within every incoming task, the 500 training originals per class are
assigned to the splits in Table~\ref{tab:setup-splits}. The assignment is
stratified and uses a fixed split seed. Clean, BrainWash and random-control
views of one original always remain together. An image cannot contribute
its clean view to regression training and its poisoned view to detector
testing.

\begin{table}[tbp]
\centering
\begin{tabular}{lrrp{6.0cm}}
\toprule
Split & Per class & Per task & Use\\
\midrule
Training & 300 & 3000 & Source regression and scaler fitting\\
Validation & 50 & 500 & Clean sample-threshold calibration on Tasks 1 and 4\\
Reserve & 100 & 1000 & Clean count-threshold calibration on Tasks 1 and 4\\
Test & 50 & 500 & Development or final detector metrics, depending on task\\
\bottomrule
\end{tabular}
\caption{Detector splits from each ten-class incoming task's 5000 training
originals. All 5000 remain available to the attack and full-task training
controls. Target training, validation and reserve partitions do not fit
the detector.}
\label{tab:setup-splits}
\end{table}

The source regression consequently receives 3000 clean negatives and 3000
BrainWash positives. These are 6000 feature rows from 3000 originals, not
6000 independently sampled images. The source and calibration attacks use
reckless BrainWash at budget 0.3. Random-control views are scored for
evaluation but are not added as training negatives. The target sample
comparison uses 500 clean and 500 poisoned detector-test views, with the
500 random-control views reported separately.

Table~\ref{tab:setup-information} records the information boundary. Avoiding
target detector fitting does not mean that every part of the experiment is
label-free: class labels are needed to train the continual learner and
construct the benchmark, and source poisoning labels train the regression.

\begin{table}[tbp]
\centering
\small
\begin{tabularx}{\textwidth}{p{5.0cm}X}
\toprule
Information & Permitted use\\
\midrule
True image classes & CL training, attack generation, stratified splits and offline accuracy; not incoming descriptor extraction\\
Source clean/poison labels & Supervised regression fitting\\
Historical clean status & Supervised validation/reserve calibration; trust assumption for inversion references\\
Target clean/poison labels & Offline mixture construction and metric calculation only\\
Incoming descriptor values & Frozen scoring or direct reference comparison; not scaler or threshold fitting\\
\bottomrule
\end{tabularx}
\caption{Information available to each part of the experiment. Historical
inversion targets are not incoming-image class labels.}
\label{tab:setup-information}
\end{table}

\subsection{Continual Learners and Pre-Update Models}
\label{subsec:setup-learners}

The backbone is the ResNet-18-style implementation supplied with the
experimental repository, based on residual networks~\cite{he2016deep}.
It is not an unmodified \texttt{torchvision} ImageNet model. Its initial
width is 32 channels, its four residual stages have widths 32, 64, 128 and
256, and adaptive $2\times2$ pooling gives a 1024-dimensional vector before
the task heads. Each task has a separate ten-class linear head. Task
identity is known, so the continual-learning evaluation is
task-incremental.

The completed comparison includes EWC~\cite{kirkpatrick2017ewc},
MAS~\cite{aljundi2018mas}, RWalk~\cite{chaudhry2018rwalk}, and the
repository's EWC-based AFEC~\cite{wang2021afec} and
ANCL~\cite{kim2023ancl} variants. Table~\ref{tab:setup-cl} gives their
registered coefficients. These are coefficients in the respective
implementations; their magnitudes are not comparable measures of
regularisation strength across methods.

\begin{table}[tbp]
\centering
\begin{tabular}{lrr}
\toprule
Implementation & Main coefficient & Auxiliary coefficient\\
\midrule
EWC & 500000 & ---\\
MAS & 10 & ---\\
RWalk & 1 & ---\\
AFEC-EWC & 500000 & 100\\
ANCL-EWC & 500000 & 100\\
\bottomrule
\end{tabular}
\caption{Coefficients fixed in the CIFAR-100 plan. AFEC-EWC and ANCL-EWC
require the implementation qualification in
Section~\ref{subsec:setup-validity}.}
\label{tab:setup-cl}
\end{table}

Victim and continuation training use 100 epochs per learned task, SGD with
learning rate 0.01, batch size 16 and gradient-norm clipping at 100.
The detector reads each checkpoint \emph{before} its incoming task is
learned. Its incoming head is reset with seed 20260720 and is held fixed
throughout extraction, including historical-reference extraction. The
feature-extraction seed is 20260721. No defender update is performed to
obtain the per-image descriptors.

The choice of 100 training epochs and the detector-specific task roles
defines this experiment. The study uses BrainWash configurations as a
controlled test bed; it does not claim an exact reproduction of every
setting or number in the original paper's main table.

\subsection{Attack Generation and Random Controls}
\label{subsec:setup-attacks}

Historical inversion generates 128 synthetic samples for each preceding
task, using 2000 optimisation iterations and the repository's
BatchNorm-regularised inversion procedure. These are 128 samples per
\emph{task}, not per class. Inversion provides historical examples for the
attack and for the reference-based detectors; the supervised regression is
trained on real source-task images instead.

For every method and seed, four target configurations combine reckless
and cautious BrainWash with $\ell_\infty$ budgets 0.1 and 0.3. These are
four configurations of one attack family, not four independent attack
families. Table~\ref{tab:setup-attack} records the common optimisation
settings. The cautious objective uses the fixed incoming-task weight
$w_{\mathrm{cur}}=1$.

\begin{table}[tbp]
\centering
\small
\begin{tabularx}{\textwidth}{p{5.5cm}X}
\toprule
Component & Setting\\
\midrule
Source/development attack & Reckless, $\epsilon=0.3$\\
Target settings & Reckless/cautious $\times$ $\epsilon\in\{0.1,0.3\}$\\
Outer attack iterations & 5000 (called epochs in the code)\\
Inner optimisation & One SGD step; learning rate 0.001\\
Attack minibatch size & 128\\
Perturbation sign-step size & 0.005, projected to the declared budget\\
Head reset during attack & Every outer iteration\\
Historical data & Inversion samples; no real historical replay in the attack\\
Saved perturbation used & \texttt{latest\_noise} with its stored image permutation\\
\bottomrule
\end{tabularx}
\caption{Attack-generation settings. The repeated attack-time head reset
is separate from the detector's single fixed-head initialisation.}
\label{tab:setup-attack}
\end{table}

The perturbation tensor is stored in the shuffled image order used by the
attacker. Extraction inverts this permutation to align each original with
its own modified view. Both detector evaluation and the poisoned
continuation use the saved latest perturbation, rather than silently
substituting the separately stored best-so-far field.

Attack optimisation sees the full incoming-task training pool, including
originals subsequently assigned to the detector test split. The test
originals are held out from \emph{detector fitting and calibration}, not
from attack generation. One perturbation artifact is generated for each
task/attack configuration and reused to construct all poisoning-fraction
simulations. The attacker is not rerun for every simulated dataset or
every fraction.

For the random control, each pixel receives independent uniform noise in
$[-\epsilon,\epsilon]$, followed by clipping the image to $[0,1]$.
The control uses the same budget as BrainWash, but its realised $\ell_2$
norm is not matched to the attack. Its extraction seed is 20260722, with
an image-index offset. The uniform-noise continuation controls use
separate draws at the same budget, not necessarily the exact random views
scored by the detector. Since random perturbations can also damage
retention, their alert rates are reported as noise sensitivity rather than
being included in clean false-positive rates.

\subsection{Supervised Fitting, Calibration and Ablations}
\label{subsec:setup-supervised}

StandardScaler and logistic regression are fitted on source-training
clean/poison rows only. Regression uses the L-BFGS solver, $C=1$, an
$\ell_2$ penalty, at most 2000 solver iterations and classifier seed
20260924. Hyperparameters are unchanged across methods and seeds; the
fitted coefficients and scaler are specific to each run.

The principal threshold rule is \texttt{history\_mad}, as defined in
Section~\ref{subsec:methods-thresholds}. It was fixed after the earlier
seed-5 pilot and before the expanded seeds 6--8. The source-quantile and
historical-quantile rules are also frozen and evaluated as comparisons.
Task 4 test data are used to check the fixed rule, not to choose a
replacement when the target result is disappointing. Final sample and
count cutoffs use full stored precision. Scores produced by different
fitted models are not assumed to share a calibrated probability scale.

Nine feature variants are fitted, including the full detector
(Table~\ref{tab:setup-ablations}). Every reduced model is retrained from
the same source originals and recalibrated using the same historical
splits. Removing a feature is not implemented by zeroing it at prediction
time. The stage variant replaces the 80 tensor norms with five stage
norms but keeps the same four context features, isolating granularity.

\begin{table}[tbp]
\centering
\small
\begin{tabularx}{\textwidth}{p{5.2cm}rX}
\toprule
Variant & Inputs & Retained information\\
\midrule
Full & 84 & 80 tensor norms and four context features\\
Gradients only & 80 & Tensor norms\\
Context only & 4 & Entropy, confidence, margin and activation norm\\
Stage norms with context & 9 & Five stage norms and four context features\\
Without entropy & 83 & Full set minus entropy\\
Without confidence & 83 & Full set minus confidence\\
Without margin & 83 & Full set minus margin\\
Without activation norm & 83 & Full set minus activation norm\\
Without uncertainty & 81 & Tensor norms and activation norm\\
\bottomrule
\end{tabularx}
\caption{Supervised feature variants. The final row removes entropy,
confidence and margin together.}
\label{tab:setup-ablations}
\end{table}

\subsection{Simulated Incoming Datasets}
\label{subsec:setup-bags}

A dataset-level evaluation samples $N=150$ distinct originals without
replacement from the 500-original detector test pool. At requested
contamination fraction $r$, the number of modified views is
\begin{equation}
    m(r)=\left\lfloor150r+\frac{1}{2}\right\rfloor.
    \label{eq:setup-rounded-count}
\end{equation}
The first $m(r)$ positions in the randomly ordered sample use the modified
view and the remaining positions use the clean view. An original therefore
appears only once in a bag, never as both a clean and a poisoned image.
Sampling is not forced to give exactly equal class counts within each bag.

\begin{table}[tbp]
\centering
\begin{tabular}{rrr}
\toprule
Requested fraction & Modified images & Realised fraction\\
\midrule
0\% & 0 & 0\%\\
1\% & 2 & 1.33\%\\
5\% & 8 & 5.33\%\\
10\% & 15 & 10\%\\
25\% & 38 & 25.33\%\\
50\% & 75 & 50\%\\
100\% & 150 & 100\%\\
\bottomrule
\end{tabular}
\caption{Requested and realised contamination for 150-image datasets. The
same counts are used for uniform-perturbation comparisons.}
\label{tab:setup-rates}
\end{table}

There are 100 repeated simulations per fraction. A zero-fraction bag is
entirely clean. Positive-fraction scenarios are constructed separately
using BrainWash or uniform-noise views. Originals recur across repetitions,
so the simulations measure conditional variability within a finite pool;
they are not 100 independent datasets or attack runs. The same clean
simulations are reused across the four attacks within a method/seed, and
repeated clean columns are not counted as additional experiments.

Rank and MMD are evaluated on identical bags within their direct
comparison. Their sampling stream is separate from the supervised
evaluation, so shared fractions alone should not be taken to imply
identical sampled membership across all three detectors.

\subsection{Reference-Based Comparison}
\label{subsec:setup-reference}

For CIFAR-100 Task 9, each of the nine historical tasks supplies 128
inversion descriptors. All reference and incoming descriptors use the
same frozen Task 9 model and the predicted-class feature protocol. Fitting
the Rank and MMD reference bundles precedes reading the incoming benchmark.
Neither bundle is calibrated using trusted clean images from the incoming
task.

Rank uses the ten descriptors and the historical-only pair-selection rule
in Section~\ref{subsubsec:methods-rank}. Its reference cap is 256 images per
historical task, so all 128 available images are retained. Each incoming
bag is tested in full. It uses 499 Gaussian multiplier draws and comparison
level 0.05. Unsupported bags are reported through coverage and are not
silently counted as accepted clean data.

MMD uses a 50\% historical preprocessing-fit fraction, a bank cap of 512
rows, 128 random Fourier components and 199 permutations. With the
1152 available historical descriptors, this leaves 576 fit rows, 512 bank
rows and 64 unused rows. The coordinate scale floor is $10^{-8}$ and the
standardised clipping limit is 12. Both tests use fixed random seeds and
the declared reference identities. The comparison is between a pooled
distribution discrepancy and taskwise rank dependence; it is not a test
of two algorithms operating on the supervised 84-feature representation.

\subsection{Evaluation Measures and Forgetting Controls}
\label{subsec:setup-metrics}

Image-level evaluation reports ROC-AUC and average precision on the
balanced clean/poison test views. Clean false-positive rate and poisoned
true-positive rate are then measured at the frozen cutoff. ROC-AUC
summarises ranking across thresholds~\cite{fawcett2006roc}; it does not
establish that the selected operating threshold transfers successfully.
Uniform-noise alert rates are reported separately.

At dataset level, the clean false-alert rate is the fraction of clean
simulations that alert. Detection at a given contamination count is the
fraction of the corresponding poisoned simulations that alert. Coverage
is reported for Rank when comparisons are unsupported. A low alert rate
cannot be interpreted as low false-positive rate unless clean and
poisoned scenarios are distinguished.

For the main study, metrics are reported per seed and summarised by their
mean and sample standard deviation across seeds 6, 7 and 8. The seed is
the repeated training/attack run; the simulated bag is not used as an
independent replicate for a population-level confidence claim. The
class order, benchmark and final task have been examined in earlier
development, so later seeds provide a robustness check within that
setting rather than independent confirmation on a new population.

Forgetting is measured by continuing the same pre-update checkpoint on
the full incoming task under clean, BrainWash and uniform-noise conditions.
Let $a_{t,j}^{(q)}$ be class-test accuracy on task $j$ after condition $q$
learns task $t$. The historical mean and backward transfer are
\begin{align}
    A_{t,\mathrm{past}}^{(q)}
       &=\frac{1}{t}\sum_{j=0}^{t-1}a_{t,j}^{(q)},\\
    \operatorname{BWT}_{t}^{(q)}
       &=\frac{1}{t}\sum_{j=0}^{t-1}
                   \bigl(a_{t,j}^{(q)}-a_{j,j}^{(q)}\bigr).
    \label{eq:setup-retention}
\end{align}
Backward transfer follows the usual continual-learning
definition~\cite{lopezpaz2017gem}. The extra historical accuracy loss is
\begin{equation}
    \Delta_t^{\mathrm{BW}}
       =A_{t,\mathrm{past}}^{(\mathrm{clean})}
        -A_{t,\mathrm{past}}^{(\mathrm{BW})}.
    \label{eq:setup-attack-damage}
\end{equation}
Differences are reported in percentage points, together with incoming-task
accuracy. These controls train on the full 5000-image incoming task, not
on the 150-image detector bags. Their measured damage is therefore not
the damage caused by, for example, eight poisoned images in a bag.
No continuation filters images using detector predictions, so these
measurements do not establish detector-based prevention of forgetting.

\subsection{Run Matrix and Reproducibility}
\label{subsec:setup-reproducibility}

The completed CIFAR-100 matrix has five implementations, three seeds and
four target attack settings: 15 method/seed cells and 60 target
configurations. Each cell has nine supervised feature variants, three
frozen threshold-rule comparisons, Rank/MMD evaluation and paired
forgetting controls. The supplied run-state records mark all 540
registered stages complete. Completion refers to execution and recorded
outputs, not to a detector meeting its intended error rate.

The archived environment uses Python 3.8.20, PyTorch 1.10.1, torchvision
0.11.2, NumPy 1.24.4, pandas 2.0.3, SciPy 1.10.1 and scikit-learn 1.3.2.
GPU work runs on Tesla V100 cards with 16 GB memory. Independent cells
can run on separate cards; a single attack is not distributed across
cards. Hardware occupancy and subprocess wall time are not reported as
GPU kernel time or as a promise of reproducibility at identical speed.

Plans record settings and source hashes. Checkpoints, inversion files,
attack permutations, feature metadata and split manifests preserve data
and model identity. Completed stages are reused only after output-hash
verification. An interrupted stage is archived and restarted; recovery
is stage-level, not an epoch checkpoint. These checks identify accidental
mixing of runs, but do not make stochastic GPU execution bitwise
deterministic.

\subsection{Full-Data ImageNet Extension}
\label{subsec:setup-extension}

The registered extension uses EWC only, two seeds (6 and 7), the same four
target attack settings and the same nine supervised feature groups.
Source fitting and clean calibration remain on Tasks 1 and 4. The final
task is 9 for miniImageNet and 19 for tinyImageNet. Feature definitions,
the form of the threshold rule and the data-split counts remain unchanged;
the scaler, regression and numerical cutoffs are fitted separately for
each dataset and seed. This tests transfer to a later task \emph{within}
each dataset. It does not test applying a CIFAR-trained detector to
ImageNet without fitting or calibration.

The plan has four dataset/seed cells and sixteen target configurations.
It retains all training images and the registered training and attack
iterations. Rank and MMD are not part of this extension. Results are to be
reported only for cells that actually finish; the plan itself provides
no evidence of detection performance or generalisation.

The new schedule avoids repeating identical clean continuation training
for all four attacks, shares uniform continuation at the same budget,
and reuses compatible clean and random-control feature measurements.
Poison features remain attack-specific, and different task checkpoints
are not interchangeable. This gives 30 stages per cell rather than the
completed CIFAR schedule's 36, without treating the reduced stage count
as a proportional runtime reduction. One GPU processes miniImageNet and
the other tinyImageNet; each processes its two seeds sequentially.

\subsection{Implementation Qualifications and Scope}
\label{subsec:setup-validity}

The detector code adapts task visibility and data loading while executing
the upstream learner and attack implementations without editing them.
This preserves the distinction between detector changes and changes to
the continual-learning baseline. It also means that limitations of an
upstream implementation remain relevant to the resulting checkpoints.

In AFEC-EWC and ANCL-EWC, source inspection identified a first-task path
that replaces the auxiliary model after its optimiser has been created,
without rebuilding that optimiser on the default path. The saved runs
must consequently be identified as results from these repository
implementations, not as independently verified reproductions of the
published auxiliary-network algorithms. Their outputs are retained and
reported with this qualification; they are not used to establish a
general claim about canonical AFEC or ANCL. The main EWC, MAS and RWalk
observations should be discussed separately from them.

Other limits concern the experimental design rather than execution.
There is one fixed class order, one attack family with four settings,
and a finite image pool reused across simulations. Poisoning labels
are available for supervised source fitting, while the reference-based
methods assume clean historical checkpoints and useful inversion
samples. No method is recalibrated on the final incoming task, no
unsupported reference comparison is converted into a clean decision,
and no filtering benefit is inferred from detection alone. These
boundaries determine which conclusions the results can support.
